\chapter{Équation du premier degré à une inconnue}


\begin{itemize}
\item Quel est le degré des équations suivantes ? \begin{itemize}
\it $(4x-1)^2=(2x+3)^2$.
\it $2x(x+5)-(x-3)^2=0$.
\it $(2x-1)^2+(x+3)^2-5(x+7)(x-7)=8$.
\it $(x+1)^3+2(x-1)^3+x^3-3x(x+1)(x-1)=0$.
\end{itemize}
\item Résoudre les équations suivantes : 
\begin{itemize}
\it $5(2x-3) - 4(5x-7)=19-2(x+11)$. 
\it $4(x+3)-7x+17 = 8(5x-1)+166$. 
\it $17-14(x+1)= 13-4(x+1)-5(x-3)$. 
\it $5x+3,5+(3x-4) = 7x-3(x-0,5)$. 
\it $7(4x+3)-4(x-1)=15(x+0,75)+7$. 
\it $17x+15(x-1)=-1-14(3x+1)$. 
\end{itemize}
\item Résoudre les équations suivantes (après développement, les termes en $x^2$ ou $x^3$ se simplifient) : 
\begin{itemize}
\it $(x-1)^2+(x+3)^2= 2(x-2)(x+1)+38$.
\it $5(x^2-2x-1)+2(3x-2)=5(x+1)^2$.
\it $(9x+1)(x-2)=(3x+4)(3x-5)$.
\it $7(3-2x)-5x(2x-1)=(5x+3)(3-2x)$. 
\it $(3x-1)^2-(2x+3)^2+7 = (2x+1)(2x-1) + x(x+7)$.
\it $(x+2)^3+(x-2)^3+(x+1)^3=3(x+1)(x+2)(x-2)$.
\end{itemize}
\item Résoudre les équations suivantes : 
\begin{itemize}
\it $\frac52x+3 - \frac{7x}4 = x + \frac94$. 
\it $\frac{3x}7 - \frac{2x}{15} + 3 = \frac{x}3+\frac{13}3$.
\it $x+\frac12-\frac{x}6 = 16 - \frac{2x}9 + \frac13$. 
\it $\frac{7x}4-2-\frac{x}2=\frac{2x}{13}-\frac{85}{52}$. 
\it $\frac{2x}3+4 -\frac{2x}5 = \frac{x}2-\frac{x}3+3,5$. 
\it $\frac{x}6-1=\frac{x}4-\frac{x}3-1$. 
\end{itemize}
\item Résoudre les équations suivantes (on se ramènera au cas entier en multipliant par un nombre opportun).
\begin{itemize}
\it $\frac{x+5}4 - \frac{x-3}6 = \frac{x}3$. 
\it $\frac{3x-7}2 + \frac{x+1}3 = -16$. 
\it $x-\frac{x+1}3 = \frac{2x+1}5$. 
\it $\frac{7-3x}{12} + \frac34 = 2(x-2) + \frac{5(5-2x)}6$. 
\it $\frac{x}5 - \frac{3x-1}6 + \frac{3-x}4 = 0$.
\it $\frac{3(x+3)}4 + \frac12 = \frac{5x+9}3 - \frac{7x-9}4$. 
\it $\frac{2x-7}5 + \frac{x+11}2 = -4$. 
\it $\frac{2x-3}3 - \frac{x-3}6 = \frac{4x+3}4 - 17$. 
\it $\frac{5x-3}4 - \frac{7x-5}9 = \frac{x+19}6$. 
\it $\frac{5x+1}8-\frac{x-1}3 = \frac{4(2x-3)}9$. 
\it $\frac{2x-1}3 - \frac{5x+2}7 = x +13$. 
\it $\frac{8x+2}5 - \frac{x-11}7 = \frac{5x-3}2 - \frac{3x-1}4$.
\it $\frac{2x-7}9 - \frac{x-5}6 = \frac{x-9}8$. 
\it $\frac{5x+7}4 - \frac{3x+5}8 = \frac{4x+9}5 - \frac{x-9}3$. 
\it $\frac{5x+6}7 - \frac{3x+1}4 = \frac{x+16}5$. 
\it $\frac{4x+7}5 -\frac{x-5}6 = \frac{2x+14}3 - \frac{2x-7}9$. 
\end{itemize}
\item Résoudre les équations suivantes (après multiplication et développement, les termes en $x^2$ disparaissent) : 
\begin{itemize}
\it $\frac{(x-1)(x+5)}3 - \frac{(x+2)(x+5)}{12} = \frac{(x-1)(x+2)}4$.
\it $\frac{(x+1)^2}3 + \frac{(x-2)(x-3)}2 = \frac{(5x-1)(x-4)}6 + \frac{28}3$. 
\it $\frac{(3x+1)(3x-1)}9 - \frac{(x-5)(x+1)}2 = \frac{(9x-1)(x+3)}{18} + \frac89$. 
\it $\frac{(4x+7)^2}4 - \frac{(5x-1)^2}7 = \frac{(8x-3)(3x+4)-79x}{56}$. 
\it $\left(x-\frac83\right)(x+0,75) = (x+4,5)(x+1,5) - \frac{145}3$. 
\it $\frac{(x-5)^2}5 + \frac{(x+3)^2}3 = \frac{(3x+1)(3x-1)-x(x+1)}{15}$. 
\it $\left(3x-\frac45\right)\left(5x+\frac23\right)
=15 (x-1)(x+1)+\frac7{15}$. 
\end{itemize}
\end{itemize}
